\section{Process Virtual Machines and Interpreted Runtimes}

The previous section followed the native execution path: a C source program is transformed into a binary executable, the operating system loads that executable into a process, and the processor executes machine instructions from the process text segment. That model is essential, but it is not the model normally used when a Python file is executed. Python introduces another execution layer between the user's source code and the physical machine.

This section explains that additional layer. When we run a Python file, the operating system does not normally load the \texttt{.py} file as a native executable. Instead, it loads the CPython interpreter executable. The user's Python source file becomes runtime input. CPython reads it, parses it, compiles it into bytecode, creates runtime execution structures, and executes that bytecode inside the interpreter process. The native process belongs to CPython; the Python program lives inside that process as bytecode, frames, objects, namespaces, and managed memory.

This architecture is an example of a \textbf{process virtual machine}. A process virtual machine is not a full computer emulator. It does not replace the operating system or boot a guest kernel. It is a software execution environment running inside an ordinary process. It defines its own virtual instruction set, its own frame model, its own object model, its own memory-management rules, and its own control-flow mechanisms. The physical CPU still executes native instructions, but those instructions belong mostly to the runtime engine.

The central representation in many such runtimes is \textbf{bytecode}: an instruction stream for a software machine rather than for the physical processor. CPython bytecode is executed by the CPython evaluation loop. JVM bytecode is loaded, verified, interpreted, and often JIT-compiled by the Java Virtual Machine. PHP source is compiled into Zend Engine opcodes. V8 uses internal bytecode-like representations and aggressively optimizes hot JavaScript paths into native code. Bash provides a contrasting model, where the shell interprets command syntax and often launches external programs as child processes.

The goal of this section is to make the boundary precise. Terms such as \emph{compiled}, \emph{interpreted}, \emph{virtual machine}, and \emph{runtime} are often used too loosely. Python is not simply ``source code executed line by line.'' CPython compiles source to bytecode and interprets that bytecode. Java is not directly executed from source; it usually runs from JVM bytecode. V8 is not merely an interpreter; it is an adaptive engine with JIT compilation. Bash is closer to command-by-command interpretation, but many commands become separate native processes.

By comparing these runtime systems, we can identify what is specific to Python and what belongs to a broader family of runtime-hosted execution models. This prepares the later chapters: Python variables are bindings to runtime objects, function calls create interpreter-managed frames, exceptions move through runtime-managed control paths, and advanced mechanisms such as generators, coroutines, and async execution exist because the runtime owns more execution state than a bare native process would expose directly.